% Generated draft: requires mathematical and source review.
\documentclass[11pt]{article}
\usepackage[margin=25mm]{geometry}
\usepackage{fontspec}
\setmainfont{Latin Modern Roman}
\usepackage{amsmath,amssymb,mathtools,unicode-math}
\setmathfont{Latin Modern Math}
\usepackage{enumitem,needspace,xurl,hyperref}
\setlength{\parindent}{0pt}
\setlength{\parskip}{5pt}
\setlength{\emergencystretch}{3em}
\newcommand{\N}{\mathbb{N}}
\newcommand{\Z}{\mathbb{Z}}
\newcommand{\Q}{\mathbb{Q}}
\newcommand{\R}{\mathbb{R}}
\newcommand{\C}{\mathbb{C}}
\newcommand{\F}{\mathbb{F}}
\newcommand{\A}{\mathbb{A}}
\newcommand{\PP}{\mathbb{P}}
\newcommand{\E}{\mathbb{E}}
\newcommand{\eps}{\varepsilon}
\newcommand{\dd}{\,\mathrm{d}}
\begin{document}
\section*{UnsolvedMath ID: 10400221}
\textbf{Source problem code:} catalogue entry
\subsection{Definitions and mathematical statement}
Let \(G\) denote the object in the source, and let \(\operatorname{Aut}(G)\) mean its structure-preserving automorphisms. Write \(\mu\) for the source's mirror map. The supplied excerpt does not define \(G\), its structure, the mirror operation, or the ambient space.

\emph{Conjecture (X.-S. Lin [262]).} Every automorphism of \(G\) is either the identity or the mirror map:
\[
 \forall\varphi\in\operatorname{Aut}(G),\qquad
 \varphi=\operatorname{id}_G\ \text{or}\ \varphi=\mu.
\]
Equivalently in the intended geometric setup, the source says that every automorphism is induced by a diffeomorphism of the ambient space. Once that space \(X\) and its action are defined, this latter assertion would say that the induced action map
\[
 \operatorname{Diff}(X)\longrightarrow\operatorname{Aut}(G)
\]
is surjective. The excerpt does not provide the hypotheses identifying the image of this map with the identity and mirror operations; that identification must be checked rather than assumed for an unspecified \(G\).

The draft retains the universal two-operation rigidity claim without replacing \(G\) by an inferred knot graph, a group, or another object. A precise definition of the ambient structures and induced action is required for a fully determined conjecture.
\subsection{Short English statement}
Lin conjectures that the only automorphisms of the source's object \(G\) are the identity and mirror map, or equivalently in its intended setup that all are induced by ambient diffeomorphisms. The object \(G\) and the ambient action are not defined in this excerpt.
\subsection{Sources}
\begin{itemize}
\item Local UnsolvedMath registry, permanent ID 10400221, source code catalogue entry; attribution X.-S. Lin [262]. The stored title is truncated after ``Conjecture 12.16 --- (X.-S.''.
\item Ohtsuki, \emph{Problems on invariants of knots and 3-manifolds} (2002), Conjecture 12.16, PDF page 166. Supplied URL: \url{https://msp.org/gtm/2002/04/gtm-2002-04-024s.pdf}; recorded access 2026-07-29; extraction pdf-font-structure. Internal reference [262] has no bibliographic entry in the record. The collection distinguishes 231 numbered bold problem, conjecture, and question headings.
\item Supplied status open and OPEN-TRIAGE. The research summary describes Lin's ambient-diffeomorphism conjecture. Current primary-literature status and redistribution terms for the publicly accessible author/publisher copy remain NEEDS\_REVIEW.
\end{itemize}
\end{document}
